\section{Conclusion}
\label{sec:conclusion}

We formulated low-resource Chinese calligraphy assessment as overlapping
document parsing followed by discrepancy-preserving registration and
human-audited comparison. Independent direction channels retain simultaneous
labels at crossings and endpoint evidence. Across matched three-seed
experiments, DeepLabV3+ leads the QC-standard split, SegFormer-B2 leads
character-disjoint direction parsing, and U-Net leads strict endpoint
localization, demonstrating that architecture choice depends on both target
and generalization protocol.

Constrained registration reduces nuisance sensitivity relative to no
alignment without implying transform invariance. The frozen score associates
modestly with blinded ratings; coverage-aware scoring is stronger, and ASDS
reaches higher retrospective internal association but remains unconfirmed on
an independent cohort. Direct-ink and parsed-union ASDS are nearly equivalent,
so the parser's value lies in directional, overlap, endpoint, and localized
evidence rather than the silhouette scalar alone. The broader
contribution is an auditable framework for limited-label visual assessment:
semantic QC, unseen-class testing, bounded registration, score-semantics
audits, human association, and explicit separation of structured evidence
from language realization.
